% Options for packages loaded elsewhere
\PassOptionsToPackage{unicode}{hyperref}
\PassOptionsToPackage{hyphens}{url}
\documentclass[
  11pt,
  letterpaper,
]{article}
\usepackage{xcolor}
\usepackage[margin=1in]{geometry}
\usepackage{amsmath,amssymb}
\setcounter{secnumdepth}{-\maxdimen} % remove section numbering
\usepackage{iftex}
\ifPDFTeX
  \usepackage[T1]{fontenc}
  \usepackage[utf8]{inputenc}
  \usepackage{textcomp} % provide euro and other symbols
\else % if luatex or xetex
  \usepackage{unicode-math} % this also loads fontspec
  \defaultfontfeatures{Scale=MatchLowercase}
  \defaultfontfeatures[\rmfamily]{Ligatures=TeX,Scale=1}
\fi
\usepackage{lmodern}
\ifPDFTeX\else
  % xetex/luatex font selection
\fi
% Use upquote if available, for straight quotes in verbatim environments
\IfFileExists{upquote.sty}{\usepackage{upquote}}{}
\IfFileExists{microtype.sty}{% use microtype if available
  \usepackage[]{microtype}
  \UseMicrotypeSet[protrusion]{basicmath} % disable protrusion for tt fonts
}{}
\makeatletter
\@ifundefined{KOMAClassName}{% if non-KOMA class
  \IfFileExists{parskip.sty}{%
    \usepackage{parskip}
  }{% else
    \setlength{\parindent}{0pt}
    \setlength{\parskip}{6pt plus 2pt minus 1pt}}
}{% if KOMA class
  \KOMAoptions{parskip=half}}
\makeatother
\usepackage{color}
\usepackage{fancyvrb}
\newcommand{\VerbBar}{|}
\newcommand{\VERB}{\Verb[commandchars=\\\{\}]}
\DefineVerbatimEnvironment{Highlighting}{Verbatim}{commandchars=\\\{\}}
% Add ',fontsize=\small' for more characters per line
\newenvironment{Shaded}{}{}
\newcommand{\AlertTok}[1]{\textcolor[rgb]{1.00,0.00,0.00}{\textbf{#1}}}
\newcommand{\AnnotationTok}[1]{\textcolor[rgb]{0.38,0.63,0.69}{\textbf{\textit{#1}}}}
\newcommand{\AttributeTok}[1]{\textcolor[rgb]{0.49,0.56,0.16}{#1}}
\newcommand{\BaseNTok}[1]{\textcolor[rgb]{0.25,0.63,0.44}{#1}}
\newcommand{\BuiltInTok}[1]{\textcolor[rgb]{0.00,0.50,0.00}{#1}}
\newcommand{\CharTok}[1]{\textcolor[rgb]{0.25,0.44,0.63}{#1}}
\newcommand{\CommentTok}[1]{\textcolor[rgb]{0.38,0.63,0.69}{\textit{#1}}}
\newcommand{\CommentVarTok}[1]{\textcolor[rgb]{0.38,0.63,0.69}{\textbf{\textit{#1}}}}
\newcommand{\ConstantTok}[1]{\textcolor[rgb]{0.53,0.00,0.00}{#1}}
\newcommand{\ControlFlowTok}[1]{\textcolor[rgb]{0.00,0.44,0.13}{\textbf{#1}}}
\newcommand{\DataTypeTok}[1]{\textcolor[rgb]{0.56,0.13,0.00}{#1}}
\newcommand{\DecValTok}[1]{\textcolor[rgb]{0.25,0.63,0.44}{#1}}
\newcommand{\DocumentationTok}[1]{\textcolor[rgb]{0.73,0.13,0.13}{\textit{#1}}}
\newcommand{\ErrorTok}[1]{\textcolor[rgb]{1.00,0.00,0.00}{\textbf{#1}}}
\newcommand{\ExtensionTok}[1]{#1}
\newcommand{\FloatTok}[1]{\textcolor[rgb]{0.25,0.63,0.44}{#1}}
\newcommand{\FunctionTok}[1]{\textcolor[rgb]{0.02,0.16,0.49}{#1}}
\newcommand{\ImportTok}[1]{\textcolor[rgb]{0.00,0.50,0.00}{\textbf{#1}}}
\newcommand{\InformationTok}[1]{\textcolor[rgb]{0.38,0.63,0.69}{\textbf{\textit{#1}}}}
\newcommand{\KeywordTok}[1]{\textcolor[rgb]{0.00,0.44,0.13}{\textbf{#1}}}
\newcommand{\NormalTok}[1]{#1}
\newcommand{\OperatorTok}[1]{\textcolor[rgb]{0.40,0.40,0.40}{#1}}
\newcommand{\OtherTok}[1]{\textcolor[rgb]{0.00,0.44,0.13}{#1}}
\newcommand{\PreprocessorTok}[1]{\textcolor[rgb]{0.74,0.48,0.00}{#1}}
\newcommand{\RegionMarkerTok}[1]{#1}
\newcommand{\SpecialCharTok}[1]{\textcolor[rgb]{0.25,0.44,0.63}{#1}}
\newcommand{\SpecialStringTok}[1]{\textcolor[rgb]{0.73,0.40,0.53}{#1}}
\newcommand{\StringTok}[1]{\textcolor[rgb]{0.25,0.44,0.63}{#1}}
\newcommand{\VariableTok}[1]{\textcolor[rgb]{0.10,0.09,0.49}{#1}}
\newcommand{\VerbatimStringTok}[1]{\textcolor[rgb]{0.25,0.44,0.63}{#1}}
\newcommand{\WarningTok}[1]{\textcolor[rgb]{0.38,0.63,0.69}{\textbf{\textit{#1}}}}
\usepackage{longtable,booktabs,array}
\usepackage{calc} % for calculating minipage widths
% Correct order of tables after \paragraph or \subparagraph
\usepackage{etoolbox}
\makeatletter
\patchcmd\longtable{\par}{\if@noskipsec\mbox{}\fi\par}{}{}
\makeatother
% Allow footnotes in longtable head/foot
\IfFileExists{footnotehyper.sty}{\usepackage{footnotehyper}}{\usepackage{footnote}}
\makesavenoteenv{longtable}
\usepackage{graphicx}
\makeatletter
\newsavebox\pandoc@box
\newcommand*\pandocbounded[1]{% scales image to fit in text height/width
  \sbox\pandoc@box{#1}%
  \Gscale@div\@tempa{\textheight}{\dimexpr\ht\pandoc@box+\dp\pandoc@box\relax}%
  \Gscale@div\@tempb{\linewidth}{\wd\pandoc@box}%
  \ifdim\@tempb\p@<\@tempa\p@\let\@tempa\@tempb\fi% select the smaller of both
  \ifdim\@tempa\p@<\p@\scalebox{\@tempa}{\usebox\pandoc@box}%
  \else\usebox{\pandoc@box}%
  \fi%
}
% Set default figure placement to htbp
\def\fps@figure{htbp}
\makeatother
\setlength{\emergencystretch}{3em} % prevent overfull lines
\providecommand{\tightlist}{%
  \setlength{\itemsep}{0pt}\setlength{\parskip}{0pt}}
\usepackage{bookmark}
\IfFileExists{xurl.sty}{\usepackage{xurl}}{} % add URL line breaks if available
\urlstyle{same}
\hypersetup{
  hidelinks,
  pdfcreator={LaTeX via pandoc}}

\author{}
\date{}

\begin{document}

\section{Emergence of Classical Constants from a Minimal Recursive
Triadic
Framework}\label{emergence-of-classical-constants-from-a-minimal-recursive-triadic-framework}

\textbf{Author:} J. W. Milton, Clarity Coalition

\textbf{Version:} 1.2

\textbf{Date:} 10 June 2026

\textbf{Keywords:} classical constants; mathematical constants;
recursive sequences; triadic recurrence; discrete dynamical systems;
Fibonacci ratio; alternating series; structural unification;
unification; convergence

\begin{center}\rule{0.5\linewidth}{0.5pt}\end{center}

\subsection{Abstract}\label{abstract}

Five classical mathematical constants (the Leibniz limit \(\pi/4\), the
golden ratio \(\varphi\), Euler's number \(e\), \(\sqrt{2}\), and
\(\ln 2\)) emerge as limits of a single family of three-variable
recurrences on a triple \((N, D, C)\) with functionally distinct roles:
a \emph{negotiation} state that is iteratively refined, a
\emph{definition/limitation} parameter that bounds, and a
\emph{contribution/integration} parameter that accumulates. These roles
are irreducible to fewer than three even when two positions share the
same numerical seed value.

Each branch of the family is specified by an initial triple
\((N_0, D_0, C_0)\) and a \emph{traversal rule} governing the evolution
of \((D_k, C_k)\). The five branches fall into three traversal classes:
Advancing (Class A, external parameter injection), Self-redefined (Class
B, internal transformation), and Fixed (Class C, constant parameters).
The limits themselves are classical: each branch reduces to a known
series, fixed point, or iterative method, so the contribution of this
framework is not in computing these numbers anew, but in
\emph{structural unification}: a single three-role grammar from which
all five mechanisms descend, combined with proved theorems (diagonal
invariance, swap symmetry, perfect-square denominators, offset
consistency, parametric families) that distinguish the five branches
from one another and from an arbitrary collection of known limits.

The \(\pi/4\) branch is structurally unique: it alone requires three
numerically distinct seeds \(\{1, 3, 5\}\) and advances its parameters
externally via a fixed step \(\Delta = 4\) rooted in a geometric axis
structure. The remaining four branches operate on seeds drawn from
\(\{0, 1, 2\}\) with purely internal or fixed parameter evolution. This
combinatorial partition is nontrivial and proved. Four open conjectures
address the classification and necessity of the admitted seeds and
traversal rules; resolving them would elevate the framework from
organizational taxonomy to a finiteness theorem within a bounded rule
class.

\begin{center}\rule{0.5\linewidth}{0.5pt}\end{center}

\subsection{1. Introduction}\label{introduction}

The Leibniz series for \(\pi/4\),

\[\frac{\pi}{4} = 1 - \frac{1}{3} + \frac{1}{5} - \frac{1}{7} + \cdots = \sum_{k=0}^{\infty} \frac{(-1)^k}{2k+1},\]

has been known since the late seventeenth century and remains a landmark
example of slowly convergent alternating series {[}Lei82, Roy11{]}.
Leibniz himself saw philosophical significance in binary arithmetic and
its connection to the I Ching {[}Lei03{]}, viewing the interplay of 0
and 1 as a model of structured emergence from minimal premises.

Five real numbers recur persistently across mathematics: \(\pi\), the
golden ratio \(\varphi = (1+\sqrt{5})/2\), Euler's number \(e\),
\(\sqrt{2}\), and \(\ln 2\). Each has multiple independent
characterizations and has been analyzed in isolation over centuries. The
question of whether a \emph{single algebraic apparatus} organizes all
five as co-equal outputs of one recursive scheme is less a question of
computing them anew than one of \emph{structural unification}: does a
common grammar exist from which they descend by varying only initial
conditions and update rules?

This paper answers that question affirmatively for the \emph{tholonic
ladder}: a family of three-variable recurrences whose branches are
distinguished by different initial triples and traversal rules, but
whose variable roles (constraining versus integrating, driven externally
versus self-contained) remain consistent across all five instances.

\textbf{What this paper provides.} Minimal definitions; a complete
branch-specification table; full proofs of convergence to each classical
limit; structural theorems (diagonal invariance, swap symmetry,
perfect-square denominators, offset consistency, parametric families)
with proofs; quantitative convergence data; and four open conjectures
that identify the remaining mathematical gap between taxonomy and
theorem.

\textbf{What this paper does not provide.} New proofs of the underlying
series identities: those are classical and referenced. The contribution
is the shared grammar and the structural contrasts.

\textbf{Organization.} §2 defines the tholonic triad and motivates its
geometry. §3 argues the irreducibility of three variables. §4 introduces
the ladder family and its traversal-class taxonomy. §5 provides the
complete branch specification and proves the five convergence
propositions. §6 establishes the structural properties. §7 records
convergence data. §8 states four open conjectures. §9 covers related
work. §10 discusses scope and implications. §11 concludes.

\begin{center}\rule{0.5\linewidth}{0.5pt}\end{center}

\subsection{2. The tholonic triad}\label{the-tholonic-triad}

\textbf{Definition 2.1} (Tholonic triad). A \emph{tholonic triad} is an
ordered triple \((N, D, C)\) of non-negative reals together with three
directed relational roles:

\begin{itemize}
\tightlist
\item
  \(N\) (\emph{negotiation}): the running state; the quantity being
  iteratively refined, emerging from the interaction of the other two.
\item
  \(D\) (\emph{definition/limitation}): the constraining, bounding
  force; what \emph{limits} the state.
\item
  \(C\) (\emph{contribution/integration}): the accumulating,
  synthesizing force; what \emph{grows} the state.
\end{itemize}

The triad is not merely a vector: the three positions have distinct
semantic roles that persist across all five branches, regardless of
whether two positions share a numerical value.

\subsubsection{2.1 Geometric origin of the axis
labels}\label{geometric-origin-of-the-axis-labels}

The tholonic geometry {[}Mil24{]} assigns successive powers of two,
\(2^0, 2^1, \ldots, 2^5\), to six vertices of a triangular
configuration: outer role vertices \(N = 2^0\), \(D = 2^1\),
\(C = 2^2\), and three inner midpoint vertices carrying
\(2^3, 2^4, 2^5\). Each directed axis runs from one outer role vertex
through an inner midpoint to another outer role vertex; its three vertex
values sum to \(7\) times a characteristic multiplier:

\begin{itemize}
\tightlist
\item
  Definition axis (\(N \to D\), through \(2^5\)):
  \(\;1 + 32 + 2 = 35 = 7 \times 5\)
\item
  Contribution axis (\(C \to N\), through \(2^4\)):
  \(\;4 + 16 + 1 = 21 = 7 \times 3\)
\item
  Instantiation axis (\(D \to C\), through \(2^3\)):
  \(\;2 + 8 + 4 = 14 = 7 \times 2\)
\end{itemize}

The axis multipliers \(5\), \(3\), and \(2\) carry the interpretation of
canonical role weights within the triad. The Instantiation axis
multiplier, denoted \(\mathrm{istep} = 2\), enters the \(\pi/4\) branch
through two derived quantities:

\[d_{\mathrm{step}} = \mathrm{istep}^2 = 4, \qquad c_{\mathrm{step}} = 2 \times \mathrm{istep} = 4.\]

Both equal \(4\), fixing a single step size \(\Delta = 4\) without free
parameters. This is not an arbitrary choice: the step arises from
squaring and doubling the same geometric constant, and it is the
smallest nontrivial step consistent with both operations yielding an
integer.

The seed triple \((N_0, D_0, C_0) = (1, 3, 5)\) of the \(\pi/4\) branch
then receives a geometric interpretation. Each axis connects two of the
three roles and excludes the third, and its multiplier supplies the seed
(or step) of precisely the excluded role: the Definition axis
(\(N \to D\)) excludes \(C\) and carries \(5 = C_0\); the Contribution
axis (\(C \to N\)) excludes \(D\) and carries \(3 = D_0\); the
Instantiation axis (\(D \to C\)) excludes \(N\) and carries
\(\mathrm{istep} = 2\), which fixes the step \(\Delta = 4\) above. The
seeds \(3\) and \(5\) are simultaneously the first members of the
\((4n-1, 4n+1)\) parity classes traversed by the recursion, and
\(N_0 = 1\) is the multiplicative identity, the generative unit from
which accumulation begins. The assignment is not merely numerical
convenience: each role draws its seed from the complementary axis, so
the axis geometry supplies both the seeds and the step that drives them.

More generally, any branch admits an interpretation in which the initial
numbers align with geometric roles, but the \(\pi/4\) branch is the only
one where the step itself is directly derivable from the axis structure
of the underlying geometry. The remaining four branches operate without
external injection; their seeds are drawn from the minimal lattice
\(\{0, 1, 2\}\) and their dynamics are either fully self-contained
(Class B) or parameter-constant (Class C). This geometric asymmetry is
part of what separates the \(\pi/4\) branch from the rest, and it
motivates the seed-partition theorem of Proposition 6.1.

\begin{figure}
\centering
\pandocbounded{\includegraphics[keepaspectratio]{figures/1_tholonic-triad.png}}
\caption{Figure 1. Left: the tholonic triad. Outer vertices carry roles
N, D, C with values \(2^0, 2^1, 2^2\); inner midpoint vertices carry
\(2^3, 2^4, 2^5\). Each directed axis sums to a multiple of 7
(multipliers 5, 3, 2). Right: the Instantiation axis multiplier istep=2
uniquely fixes the pi/4 branch step delta=4 and the seed triple (N\_0,
D\_0, C\_0) = (1, 3, 5).}
\end{figure}

\begin{center}\rule{0.5\linewidth}{0.5pt}\end{center}

\subsection{3. The irreducible three-variable
structure}\label{the-irreducible-three-variable-structure}

\textbf{Lemma 3.1} (Triadic irreducibility). \emph{Let \(\mathcal{R}\)
be a recurrence on \(m\) real variables \(x^{(1)}, \ldots, x^{(m)}\) in
which the update to one distinguished state variable \(x^{(1)}\) is of
the form}

\[x^{(1)}_{k+1} = g\!\left(x^{(1)}_k;\, \alpha_k, \beta_k\right)\]

\emph{where \(\alpha_k\) limits the correction magnitude and \(\beta_k\)
controls the additive or multiplicative integration of each step.
Suppose that (i) \(\alpha_k\) and \(\beta_k\) are functionally
independent in the sense that \(g\) is not expressible as a function of
a single combined argument \(h(\alpha_k, \beta_k)\) for all admissible
\((\alpha_k, \beta_k)\), and (ii) the named limit
\(L = \lim_{k\to\infty} x^{(1)}_k\) is not expressible as a trivial
fixed point of a one- or two-variable recurrence. Then \(\mathcal{R}\)
requires \(m \geq 3\): a state variable and two functionally independent
auxiliary variables.}

\emph{Proof.} The distinguished variable \(x^{(1)}\) must persist across
iterations to record the refined state, so it cannot be eliminated.
Suppose \(m = 2\). Then both control arguments are determined by the
single auxiliary variable \(y_k\): \(\alpha_k = a(y_k)\) and
\(\beta_k = b(y_k)\), and the admissible pairs \((\alpha_k, \beta_k)\)
are confined to the one-parameter image of \(y \mapsto (a(y), b(y))\).
On that set, \(g\) factors through the single argument \(y\) itself:
\(g(x;\, \alpha_k, \beta_k) = \tilde{g}(x;\, y_k)\), exhibiting exactly
the combined-argument factorization that hypothesis (i) excludes on the
admissible set. Hypothesis (ii) rules out the degenerate alternative in
which the controls are constant or absent and \(L\) is reachable with
fewer variables. Hence two functionally independent auxiliary variables
are required, and with the state variable, \(m \geq 3\). \(\square\)

A critical structural property of the tholonic framework is that
\textbf{every branch requires exactly three variables}, regardless of
whether some share the same numerical value.

In the \(e\) branch, seeds are \(N_0 = 0\), \(D_0 = 1\), \(C_0 = 1\);
superficially reducible to two values. In the \(\sqrt{2}\) branch,
\(D_0 = C_0 = 2\); apparently a single parameter. However, the triad is
not about three \emph{numbers}; it is about three \emph{functional
roles}:

\begin{enumerate}
\def\labelenumi{\arabic{enumi}.}
\tightlist
\item
  \textbf{\(N\)}: the state being negotiated; what emerges.
\item
  \textbf{\(D\)}: the defining boundary; what constrains.
\item
  \textbf{\(C\)}: the integrating accumulator; what synthesizes.
\end{enumerate}

In the \(e\) branch, \(D = 1\) is the fixed numerator that never
changes; the unchanging boundary. \(C\) starts at 1 but grows as a
factorial; it is the expanding denominator that absorbs structure each
iteration. They begin at the same value but diverge immediately in
behavior: one constrains, the other integrates. Collapsing them into a
single variable would destroy the distinction between the limiting and
contributing functions: the formula would still yield the same number,
but the structural reason the recursion works would be erased.

In the \(\sqrt{2}\) branch, \(D = 2\) is the target value being bounded
toward; the defining constraint. \(C = 2\) is the averaging divisor; the
integrator that synthesizes the overshoot and undershoot. Without
\(D/N\), there is no correction. Without division by \(C\), there is no
synthesis. The contribution/integration role is not always addition; in
\(\sqrt{2}\) it averages, in \(\pi/4\) it adds, but both are acts of
integration.

\textbf{The triad is irreducible: a state, a limiter, and a contributor
constitute the minimum structure for a recursion that converges to a
non-trivial constant.} Two variables give either unconstrained growth or
static limitation; three give the dynamic tension that converges.

\begin{center}\rule{0.5\linewidth}{0.5pt}\end{center}

\subsection{4. The ladder family}\label{the-ladder-family}

\textbf{Definition 4.1} (Tholonic ladder recurrence). Given
\((N_0, D_0, C_0) \in \mathbb{R}^3\), a \emph{tholonic ladder branch} is
a discrete dynamical system

\[N_{k+1} = f(N_k;\, D_k,\, C_k), \qquad k \in \mathbb{N}_0,\]

together with a \emph{traversal rule} specifying how \((D_k, C_k)\) is
updated at each step.

\textbf{Definition 4.2} (Branch). A \emph{branch} is the tuple
\(\bigl((N_0, D_0, C_0),\, f,\, \text{traversal rule}\bigr)\).

\subsubsection{4.1 Classification of traversal
types}\label{classification-of-traversal-types}

The five documented branches fall into three classes:

\textbf{Class A: Advancing.} \((D_k, C_k)\) each receive a fixed
external increment at every step. New information is injected per
iteration; the recursion is not self-contained. Only the \(\pi/4\)
branch belongs to this class.

\textbf{Class B: Self-redefined.} \((D_k, C_k)\) evolve by a
transformation of their own current values (Fibonacci swap for
\(\varphi\), factorial scaling for \(e\)) but receive no external
increment. The factorial rule consumes the iteration counter \(k\) as a
multiplier, but injects no new constant beyond it; the recursion
otherwise feeds on its own previous state.

\textbf{Class C: Fixed.} \((D_k, C_k)\) are held constant at their seed
values for all \(k\). Convergence is driven by the update map \(f\)
alone: as a pure fixed-point iteration in the \(\sqrt{2}\) branch (where
\(k\) never enters \(f\)), or with the iteration index entering \(f\)
directly in the \(\ln 2\) branch.

This trichotomy partitions the five branches cleanly: one Class A
(\(\pi/4\)), two Class B (\(\varphi\), \(e\)), two Class C
(\(\sqrt{2}\), \(\ln 2\)). The classification can equivalently be stated
on two independent axes: whether \((D_k, C_k)\) evolve at all, and
whether the branch consumes the iteration counter \(k\).

\begin{longtable}[]{@{}llll@{}}
\toprule\noalign{}
Branch & \((D, C)\) evolve? & Counter \(k\) consumed? & Class \\
\midrule\noalign{}
\endhead
\bottomrule\noalign{}
\endlastfoot
\(\pi/4\) & yes (external increment \(\Delta\)) & no & A \\
\(\varphi\) & yes (internal swap) & no & B \\
\(e\) & yes (scaling by \(k\)) & yes, in traversal & B \\
\(\sqrt{2}\) & no & no & C \\
\(\ln 2\) & no & yes, in \(f\) & C \\
\end{longtable}

Class A is evolution by external constant; Class B is evolution by
internal or counter-indexed transformation; Class C is no evolution.

\begin{figure}
\centering
\pandocbounded{\includegraphics[keepaspectratio]{figures/1_traversal-classes.png}}
\caption{Figure 2. Classification of the five tholonic ladder branches
by traversal type. Class A (red) injects new information each step via
an external increment. Class B (blue) evolves parameters by an internal
rule (Fibonacci swap or factorial growth). Class C (green) holds all
parameters fixed; convergence is driven by the update map alone, with
the iteration index entering f directly in the ln 2 branch.}
\end{figure}

\begin{center}\rule{0.5\linewidth}{0.5pt}\end{center}

\subsection{5. Branch specification and convergence
proofs}\label{branch-specification-and-convergence-proofs}

\textbf{Table 1.} Complete branch specification.

\begin{longtable}[]{@{}
  >{\raggedright\arraybackslash}p{(\linewidth - 12\tabcolsep) * \real{0.1129}}
  >{\raggedright\arraybackslash}p{(\linewidth - 12\tabcolsep) * \real{0.1129}}
  >{\raggedright\arraybackslash}p{(\linewidth - 12\tabcolsep) * \real{0.1129}}
  >{\raggedright\arraybackslash}p{(\linewidth - 12\tabcolsep) * \real{0.1129}}
  >{\raggedright\arraybackslash}p{(\linewidth - 12\tabcolsep) * \real{0.1774}}
  >{\raggedright\arraybackslash}p{(\linewidth - 12\tabcolsep) * \real{0.2581}}
  >{\raggedright\arraybackslash}p{(\linewidth - 12\tabcolsep) * \real{0.1129}}@{}}
\toprule\noalign{}
\begin{minipage}[b]{\linewidth}\raggedright
Limit
\end{minipage} & \begin{minipage}[b]{\linewidth}\raggedright
\(N_0\)
\end{minipage} & \begin{minipage}[b]{\linewidth}\raggedright
\(D_0\)
\end{minipage} & \begin{minipage}[b]{\linewidth}\raggedright
\(C_0\)
\end{minipage} & \begin{minipage}[b]{\linewidth}\raggedright
\(N_{k+1}\)
\end{minipage} & \begin{minipage}[b]{\linewidth}\raggedright
Traversal rule
\end{minipage} & \begin{minipage}[b]{\linewidth}\raggedright
Class
\end{minipage} \\
\midrule\noalign{}
\endhead
\bottomrule\noalign{}
\endlastfoot
\(\pi/4\) & \(1\) & \(3\) & \(5\) & \(N - \dfrac{1}{D} + \dfrac{1}{C}\)
& \(D \leftarrow D+4,\quad C \leftarrow C+4\) & A \\
\(\varphi\) & \(1\) & \(1\) & \(2\) & \(N_0 + \dfrac{D}{C}\) &
\(D \leftarrow C,\quad C \leftarrow C + D\) & B \\
\(e\) & \(0\) & \(1\) & \(1\) & \(N + \dfrac{D}{C}\) &
\(C \leftarrow C \cdot \max(k,1)\) & B \\
\(\sqrt{2}\) & \(1\) & \(2\) & \(2\) & \(\dfrac{N + D/N}{C}\) & \(D, C\)
fixed & C \\
\(\ln 2\) & \(0\) & \(1\) & \(1\) & \(N + (-1)^k \dfrac{D}{k + C}\) &
\(D, C\) fixed & C \\
\end{longtable}

In every branch \(D\) plays the bounding/constraining role and \(C\) the
integrating/synthesizing role, despite the operations differing
(subtraction in \(\pi/4\), averaging in \(\sqrt{2}\), factorial growth
in \(e\), Fibonacci swap in \(\varphi\), harmonic stepping in
\(\ln 2\)). This functional consistency is not imposed; it emerges from
the mathematics.

\subsubsection{\texorpdfstring{5.1 The \(\pi/4\)
branch}{5.1 The \textbackslash pi/4 branch}}\label{the-pi4-branch}

\textbf{Proposition 5.1.} \emph{With \((N_0, D_0, C_0) = (1, 3, 5)\) and
step \(\Delta = 4\),}

\[N_\infty = 1 + \sum_{n=1}^{\infty}\!\left(\frac{1}{4n+1} - \frac{1}{4n-1}\right) = 1 + \sum_{n=1}^{\infty}\frac{-2}{16n^2-1} = \frac{\pi}{4}.\]

\emph{Proof.} At iteration \(n \geq 1\), \(D_n = 4n-1\) and
\(C_n = 4n+1\). The update adds \(-1/(4n-1) + 1/(4n+1)\) to the running
total, which equals \(-2/[(4n-1)(4n+1)] = -2/(16n^2-1)\). Starting from
\(N_0 = 1 = 1/1 - (-1/3) + \cdots\) (the first Leibniz term), the
accumulated sum is a grouping of consecutive Leibniz pairs. By the
alternating series theorem applied to
\(\sum_{k=0}^\infty (-1)^k/(2k+1)\), the grouped sum converges to
\(\pi/4\) {[}Kno56{]}. \(\square\)

\textbf{Remark} (Primitive invariant). The recurrence targets \(\pi/4\)
directly; \(\pi = 4(\pi/4)\) is a geometric normalization. We treat
\(\pi/4\) as the \emph{primitive invariant} of this branch.

\textbf{Corollary 5.2} (Perfect-square form). Multiplying both sides of
Proposition 5.1 by \(8\):

\[\pi = \sum_{n=1}^{\infty} \frac{8}{16n^2 - 1},\]

expressing \(\pi\) in terms of perfect squares offset by unity. See also
Proposition 6.4.

\textbf{Functional roles.} \(D_k\) is the odd denominator that
\emph{subtracts} from \(N_k\), constraining the total; \(C_k\) is the
successive odd denominator that \emph{adds}, contributing to it. The two
parameters advance together in lockstep, driven by the external step
\(\Delta = 4\).

\begin{figure}
\centering
\pandocbounded{\includegraphics[keepaspectratio]{figures/1_pi4-convergence.png}}
\caption{Figure 3. Partial sums \(N_k\) of the pi/4 branch converging to
pi/4 (dashed). The damped alternating behavior is characteristic of
conditionally convergent series; amplitude decays as O(1/k). Red
diamonds mark k in \{5, 25, 125\}; the shaded band shows the ±1/(2k+1)
envelope.}
\end{figure}

\subsubsection{\texorpdfstring{5.2 The \(\varphi\)
branch}{5.2 The \textbackslash varphi branch}}\label{the-varphi-branch}

\textbf{Proposition 5.3.} \emph{With \((N_0, D_0, C_0) = (1, 1, 2)\),
update \(N_{k+1} = N_0 + D_k/C_k\), and the Fibonacci traversal
\(D \leftarrow C\), \(C \leftarrow C + D\), the limit is
\(N_\infty = \varphi = (1+\sqrt{5})/2\).}

\emph{Proof.} With seeds \((D_0, C_0) = (1, 2) = (F_2, F_3)\), the
traversal maintains \((D_k, C_k) = (F_{k+2}, F_{k+3})\), consecutive
Fibonacci numbers. Ratios of consecutive Fibonacci numbers converge:
\(F_{k+2}/F_{k+3} \to 1/\varphi\) {[}HW79, Kos01{]}, with error
contracting geometrically by the factor \(\varphi^{-2} \approx 0.382\)
per step (Binet's formula). Hence
\(N_{k+1} = 1 + F_{k+2}/F_{k+3} \to 1 + 1/\varphi = \varphi\), by the
defining identity \(\varphi = 1 + 1/\varphi\). \(\square\)

\textbf{Functional roles.} \(D\) is the smaller (prior) Fibonacci term;
the boundary, the definition. \(C\) is the larger accumulated term; the
integration of the two most recent generations. Though both are
reassigned each iteration, \(D\) always carries the constraining role
and \(C\) the integrating role. The update consumes the ratio \(D/C\)
directly, so the Fibonacci traversal is the engine of convergence, not
an accessory to it.

\subsubsection{\texorpdfstring{5.3 The \(\sqrt{2}\)
branch}{5.3 The \textbackslash sqrt\{2\} branch}}\label{the-sqrt2-branch}

\textbf{Proposition 5.4.} \emph{With \((N_0, D_0, C_0) = (1, 2, 2)\) and
fixed parameters, \(N_\infty = \sqrt{2}\), with quadratic convergence.}

\emph{Proof.} The update \(N_{k+1} = (N_k + 2/N_k)/2\) is the
Newton--Babylonian (Heron) method for \(x^2 = 2\) {[}Hea21, BF15{]}. Let
\(g(x) = (x + 2/x)/2\). Then \(g(\sqrt{2}) = \sqrt{2}\) and
\(g'(\sqrt{2}) = (1 - 2/x^2)/2\big|_{x=\sqrt{2}} = 0\), confirming
quadratic convergence. \(\square\)

\textbf{Functional roles.} \(D = 2\) is the target being bounded toward;
the defining constraint. \(C = 2\) is the averaging divisor; it
synthesizes the overshoot and undershoot of the current estimate. Though
\(D = C\) numerically, their operational roles are distinct: removing
\(D/N\) eliminates the corrective force; removing the \(C\)-division
eliminates the synthesis.

\subsubsection{\texorpdfstring{5.4 The \(e\)
branch}{5.4 The e branch}}\label{the-e-branch}

\textbf{Proposition 5.5.} \emph{With \((N_0, D_0, C_0) = (0, 1, 1)\) and
\(C_k = k!\), \(N_\infty = e\).}

\emph{Proof.} The update \(N_{k+1} = N_k + D/C_k = N_k + 1/k!\)
accumulates \(\sum_{k=0}^{K} 1/k!\). As \(K \to \infty\), this sum
converges to \(e\) {[}Rud76{]}. \(\square\)

\textbf{Functional roles.} \(D = 1\) is the fixed numerator; the
unchanging boundary, the definition that never varies. \(C\) starts at
\(1\) but grows as a factorial: the expanding denominator that absorbs
more structure each iteration, increasingly limiting how much of the
unit contributes. They start at the same value but diverge immediately
in behavior: one constrains, the other integrates.

\textbf{Remark} (Offset variant). When initialized with \(N_0 = 2\) (the
Instantiation axis value in the tholonic geometry), the branch converges
to \(e + 2\). The offset \(+2\) equals \(N_0\) exactly (Proposition
6.5), confirming that the raw output encodes both the target constant
and a baseline contribution from the initial triad state. In the
canonical form with \(N_0 = 0\), no offset appears.

\subsubsection{\texorpdfstring{5.5 The \(\ln 2\)
branch}{5.5 The \textbackslash ln 2 branch}}\label{the-ln-2-branch}

\textbf{Proposition 5.6.} \emph{With \((N_0, D_0, C_0) = (0, 1, 1)\) and
fixed parameters, \(N_\infty = \ln 2\).}

\emph{Proof.} The update accumulates \((-1)^k \cdot 1/(k+1)\) starting
at \(k = 0\), giving the alternating harmonic series
\(\sum_{k=0}^\infty (-1)^k/(k+1) = \sum_{j=1}^\infty (-1)^{j+1}/j = \ln 2\)
{[}Kno56{]}. \(\square\)

\textbf{Functional roles.} \(D = 1\) is the fixed numerator; the
defining boundary. \(C = 1\) is the offset in the denominator; the
integration base that steps through the harmonic denominators. They are
numerically identical but structurally distinct: one limits, one
integrates.

\textbf{Remark} (Offset variant). When \(N_0 = 1/2\), the branch
converges to \(\ln 2 + 1/2\). The offset equals \(N_0\) exactly, as it
must: the update is pure accumulation, so the seed enters the limit
additively (Proposition 6.5).

\begin{center}\rule{0.5\linewidth}{0.5pt}\end{center}

\subsection{6. Structural properties}\label{structural-properties}

The following propositions establish algebraic structure that
distinguishes the ladder family from an arbitrary collection of
convergent recurrences.

\textbf{Proposition 6.1} (Seed partition). \emph{Among the five
branches, the \(\pi/4\) branch is the unique branch with three
numerically distinct seeds; its seed set is \(\{1, 3, 5\}\). The
remaining four branches each have seeds drawn from \(\{0, 1, 2\}\).}

\begin{longtable}[]{@{}llll@{}}
\toprule\noalign{}
Branch & Seeds & Distinct values & Seed set \\
\midrule\noalign{}
\endhead
\bottomrule\noalign{}
\endlastfoot
\(\pi/4\) & \((1,3,5)\) & 3 & \(\{1,3,5\}\) \\
\(\varphi\) & \((1,1,2)\) & 2 & \(\{1,2\}\) \\
\(e\) & \((0,1,1)\) & 2 & \(\{0,1\}\) \\
\(\sqrt{2}\) & \((1,2,2)\) & 2 & \(\{1,2\}\) \\
\(\ln 2\) & \((0,1,1)\) & 2 & \(\{0,1\}\) \\
\end{longtable}

The value \(1\) is the most common seed, appearing 10 out of 15 times
across the table. The four non-\(\pi/4\) constants can all be generated
from the most primitive inputs: \(0\) (nothingness), \(1\) (unity), and
\(2\) (first duality). Only \(\pi/4\) demands the richer seed set
\(\{1, 3, 5\}\) and external parameter injection.

\textbf{Remark on seeds and primality.} The seeds \(1, 3, 5\) of the
\(\pi/4\) branch are the first three odd integers \(\geq 1\);
equivalently, the first three values in the generative sequence from
which all subsequent odd integers derive. Whether \(1\) is treated as
``prime'' or not in the analytic sense does not affect the convergence
proof; it is distinguished here as the generative unit (multiplicative
identity, its own reciprocal, its own square) from which the triadic
recursion begins.

\textbf{Proposition 6.2} (Diagonal invariance). \emph{For the Class A
recurrence template with \(D_k = C_k\) for all \(k\), the partial sums
telescope to zero net change and \(N_\infty = N_0\). Structural
non-triviality requires \(D \neq C\).}

\emph{Proof.} At each step the contribution is \(-1/D_k + 1/C_k = 0\)
when \(D_k = C_k\). The running total is unchanged for all \(k\).
\(\square\)

\textbf{Proposition 6.3} (Swap symmetry). \emph{For the Class A
recurrence with fixed \(N_0\), let \(N_\infty(D_0, C_0)\) denote the
limit with seeds \((D_0, C_0)\). Then}

\[N_\infty(D_0, C_0) + N_\infty(C_0, D_0) = 2\,N_0.\]

\emph{Proof.} Swapping \(D \leftrightarrow C\) negates the signed
contribution \(-1/D_k + 1/C_k \to -1/C_k + 1/D_k\) at every step, so the
two accumulated partial sums are antisymmetric around \(N_0\) at every
finite \(k\). The identity holds in the limit. \(\square\)

\begin{figure}
\centering
\pandocbounded{\includegraphics[keepaspectratio]{figures/1_swap-symmetry.png}}
\caption{Figure 4. Swap symmetry (Proposition 6.3): the two runs
(D\_0,C\_0)=(3,5) and (5,3) mirror around the midline N\_0=1 at every
iteration k. Run A decreases monotonically to pi/4; Run B increases
monotonically to 2-pi/4. Equal brackets confirm N\_A(k)+N\_B(k)=2N\_0=2
at each step. Inset: numerical residual is at floating-point zero
throughout.}
\end{figure}

\textbf{Proposition 6.4} (Perfect-square denominators). \emph{In the
\(\pi/4\) branch, the product \(D_n \cdot C_n\) at step \(n \geq 1\)
satisfies}

\[D_n \cdot C_n = (4n-1)(4n+1) = 16n^2 - 1.\]

\emph{This yields the closed-form identity
\(\pi = \displaystyle\sum_{n=1}^{\infty} \dfrac{8}{16n^2 - 1}\),
expressing \(\pi\) in terms of perfect squares offset by unity.}

\emph{Proof.} With \(D_n = 3 + 4(n-1) = 4n-1\) and
\(C_n = 5 + 4(n-1) = 4n+1\), the product is \((4n-1)(4n+1) = 16n^2-1\).
The closed form for \(\pi\) follows directly from Proposition 5.1 by
algebraic rearrangement. \(\square\)

\begin{longtable}[]{@{}lllll@{}}
\toprule\noalign{}
\(n\) & \(D_n\) & \(C_n\) & \(D_n \cdot C_n\) & Perfect-square form \\
\midrule\noalign{}
\endhead
\bottomrule\noalign{}
\endlastfoot
1 & 3 & 5 & 15 & \(16(1)^2 - 1\) \\
2 & 7 & 9 & 63 & \(16(2)^2 - 1\) \\
3 & 11 & 13 & 143 & \(16(3)^2 - 1\) \\
\end{longtable}

\begin{figure}
\centering
\pandocbounded{\includegraphics[keepaspectratio]{figures/1_perfect-squares.png}}
\caption{Figure 5. The denominator products D\_n * C\_n = (4n-1)(4n+1) =
16n\^{}2 - 1 of the pi/4 branch at scales K = 10, 100, 1000 (left to
right). Each scale reproduces the same parabolic shape, reflecting the
self-similar structure of the recursion. The dotted red curve shows the
reference parabola 16n\^{}2; the products are uniformly offset by -1
(perfect squares minus unity).}
\end{figure}

\textbf{Proposition 6.5} (Offset consistency, translation equivariance).
\emph{Let a branch have a pure-accumulation update
\(N_{k+1} = N_k + t_k\), where the term \(t_k\) depends only on
\((D_k, C_k, k)\) and not on \(N_k\) (the \(e\) and \(\ln 2\) branches).
Then the limit is translation-equivariant in the seed:}

\[N_\infty(N_0) = N_0 + N_\infty(0),\]

\emph{so the affine offset of \(N_\infty\) from the target constant
equals \(N_0\) exactly.}

\emph{Proof.} By induction, \(N_K = N_0 + \sum_{k=0}^{K-1} t_k\), and
the sum is independent of \(N_0\). Letting \(K \to \infty\) gives the
identity. \(\square\)

\begin{longtable}[]{@{}llll@{}}
\toprule\noalign{}
Branch & Variant \(N_0\) & \(N_\infty\) & Offset \\
\midrule\noalign{}
\endhead
\bottomrule\noalign{}
\endlastfoot
\(e\) & \(2\) & \(e + 2\) & \(+2 = N_0\) \\
\(\ln 2\) & \(1/2\) & \(\ln 2 + 1/2\) & \(+1/2 = N_0\) \\
\end{longtable}

In both cases the offset is the initial baseline of the accumulation,
not a failure of the framework. In the canonical forms of Table 1 with
\(N_0 = 0\), both converge directly to their targets.

\textbf{Proposition 6.6} (Parametric families). \emph{Two branches embed
in one-parameter families with the same traversal class and role
structure:}

\emph{(i) Replacing \(D = 2\) by \(D = a > 0\) in the \(\sqrt{2}\)
branch yields \(N_\infty = \sqrt{a}\), with quadratic convergence, for
every \(a\).}

\emph{(ii) Replacing the additive base \(N_0 = 1\) by
\(m \in \mathbb{N}\) in the \(\varphi\) branch, with the generalized
traversal \(D \leftarrow C\), \(C \leftarrow mC + D\), yields the
metallic means \(N_\infty = \sigma_m = (m + \sqrt{m^2 + 4})/2\); the
case \(m = 1\) recovers \(\varphi\).}

\emph{Proof.} (i) The update \(N_{k+1} = (N_k + a/N_k)/C\) with
\(C = 2\) is Newton's method for \(x^2 = a\) {[}BF15{]}; the argument of
Proposition 5.4 applies verbatim, with \(g'(\sqrt{a}) = 0\) giving
quadratic convergence. (ii) The traversal maintains consecutive terms of
the generalized Fibonacci recurrence \(G_{k+1} = m G_k + G_{k-1}\),
whose consecutive ratios converge (for any positive seeds, since
\(\sigma_m\) is the dominant root) to the positive root of
\(x^2 = mx + 1\): thus \(D_k / C_k \to 1/\sigma_m\), and
\(N_{k+1} = m + D_k/C_k \to m + 1/\sigma_m = \sigma_m\) by the defining
identity \(\sigma_m = m + 1/\sigma_m\). \(\square\)

These deformations show that the documented branches are distinguished
points inside parametric families rather than isolated curiosities, and
they make precise why the local-isolation reformulation of Conjecture
8.3 is the natural one: within a family, varying the parameter moves the
limit continuously through other algebraic targets, so uniqueness claims
must be made relative to fixed parameters.

\begin{center}\rule{0.5\linewidth}{0.5pt}\end{center}

\subsection{7. Convergence}\label{convergence}

\textbf{Table 2.} Numerical convergence at \(K = 10^5\) iterations.

\begin{longtable}[]{@{}
  >{\raggedright\arraybackslash}p{(\linewidth - 8\tabcolsep) * \real{0.1600}}
  >{\raggedright\arraybackslash}p{(\linewidth - 8\tabcolsep) * \real{0.3600}}
  >{\raggedright\arraybackslash}p{(\linewidth - 8\tabcolsep) * \real{0.1600}}
  >{\raggedright\arraybackslash}p{(\linewidth - 8\tabcolsep) * \real{0.2000}}
  >{\raggedright\arraybackslash}p{(\linewidth - 8\tabcolsep) * \real{0.1200}}@{}}
\toprule\noalign{}
\begin{minipage}[b]{\linewidth}\raggedright
Branch
\end{minipage} & \begin{minipage}[b]{\linewidth}\raggedright
\(N_K\) (computed)
\end{minipage} & \begin{minipage}[b]{\linewidth}\raggedright
Target
\end{minipage} & \begin{minipage}[b]{\linewidth}\raggedright
Residual
\end{minipage} & \begin{minipage}[b]{\linewidth}\raggedright
Rate
\end{minipage} \\
\midrule\noalign{}
\endhead
\bottomrule\noalign{}
\endlastfoot
\(\pi/4\) & \(0.785\,399\,4\ldots\) & \(0.785\,398\,2\ldots\) &
\(1.2 \times 10^{-6}\) & \(O(1/K)\) \\
\(\varphi\) & \(1.618\,033\,9\ldots\) & \(1.618\,033\,9\ldots\) &
\(< 10^{-15}\) & linear, rate \(\varphi^{-2}\) \\
\(e\) & \(2.718\,281\,8\ldots\) & \(2.718\,281\,8\ldots\) &
\(< 10^{-15}\) & \(O(1/K!)\) \\
\(\sqrt{2}\) & \(1.414\,213\,6\ldots\) & \(1.414\,213\,6\ldots\) &
\(< 10^{-15}\) & quadratic \\
\(\ln 2\) & \(0.693\,142\,2\ldots\) & \(0.693\,147\,2\ldots\) &
\(5.0 \times 10^{-6}\) & \(O(1/K)\) \\
\end{longtable}

The slow convergence of \(\pi/4\) and \(\ln 2\) (\(O(1/k)\) term decay)
is intrinsic to their alternating-series structure. The three remaining
branches converge to machine precision well within \(10^5\) iterations.
Series-acceleration methods (Euler, van Wijngaarden, Levin) would reduce
residuals for the two slow branches but are not required by the
framework.

\textbf{Explicit error bounds.} For the two alternating branches the
alternating series theorem supplies a priori envelopes: after \(k\)
iterations, \(|N_k - \pi/4| \leq 1/(4k+3)\) (the first omitted Leibniz
term) and \(|N_k - \ln 2| \leq 1/(k+1)\). At \(K = 10^5\) these give
\(2.5 \times 10^{-6}\) and \(10^{-5}\), consistent with the observed
residuals \(1.2 \times 10^{-6}\) and \(5.0 \times 10^{-6}\) (roughly
half the envelope, as expected for an alternating tail). For
\(\sqrt{2}\), Newton's method gives the exact one-step relation
\(|N_{k+1} - \sqrt{2}| = |N_k - \sqrt{2}|^2 / (2N_k)\); for \(\varphi\),
the error contracts geometrically by the factor
\(\varphi^{-2} \approx 0.382\) per step.

\begin{figure}
\centering
\pandocbounded{\includegraphics[keepaspectratio]{figures/1_convergence-comparison.png}}
\caption{Figure 6. Convergence of all five branches: log10\textbar N\_k
- N\_inf\textbar{} vs k at K=10\^{}5 iterations. Class B branches (e,
phi) reach machine precision within tens of steps; Class C
Newton-Babylonian (sqrt2) converges quadratically and hits the floor by
k\textasciitilde6. Class A (pi/4) and the Class C alternating branch (ln
2) decay as O(1/k), tracking the dotted slope-(-1) reference line
throughout.}
\end{figure}

\textbf{Convergence and traversal class.} The two slowest branches are
the Class A branch (\(\pi/4\)) and one Class C branch (\(\ln 2\)). The
fastest branches are Class B (\(e\), super-exponential; \(\varphi\),
linear with geometric ratio \(\varphi^{-2} \approx 0.382\)) and the
other Class C branch (\(\sqrt{2}\), quadratic). Whether traversal class
systematically predicts convergence order is a regularity worth
formalizing as future work.

\begin{center}\rule{0.5\linewidth}{0.5pt}\end{center}

\subsection{8. Open problems}\label{open-problems}

The structural results of §6 are proved. The following questions remain
open; positive resolutions would move the framework from taxonomy to
theorem.

\textbf{Conjecture 8.1} (Uniqueness of the \(\pi/4\) seed). \emph{The
seed triple \((1, 3, 5)\) is the unique Class A admissible seed with
\(D_0 < C_0\), \(C_0 - D_0 = 2\), \(D_0\) and \(C_0\) consecutive
members of the \(\{4n-1\}\) and \(\{4n+1\}\) sequences starting from
\(n=1\), that yields \(N_\infty = \pi/4\) under step size
\(\Delta = 4\).}

\textbf{Conjecture 8.2} (Exclusion of \(\pi/4\) from the \(\{0,1,2\}\)
lattice). \emph{No triple \((N_0, D_0, C_0) \in \{0,1,2\}^3\) with any
Class B or Class C traversal rule drawn from the admitted vocabulary
yields \(N_\infty = \pi/4\). That is, \(\pi/4\) is genuinely
inaccessible from the primitive seed lattice.}

\textbf{Conjecture 8.3} (Finite classification of named limits).
\emph{There exists a finite rule set \(\mathcal{R}\) (specifying the
allowed forms of \(f\) and the admitted traversal types over a bounded
seed space) such that the only convergent branches with limits in
\(\{\pi/4, \varphi, e, \sqrt{2}, \ln 2\}\) are exactly the five
documented here.}

If Conjecture 8.3 fails, the appropriate reformulation is: bound the
parameter freedom and show the five branches are \emph{locally isolated}
within the admissible space; no continuous deformation of seeds and rule
parameters moves one branch into another while preserving its limit.

\textbf{Conjecture 8.4} (Non-circularity of the \(\pi/4\) grouping).
\emph{The pairing \(N_{k+1} = N_k - 1/D_k + 1/C_k\) with Class A
advancement can be derived from the constraints (a) \(f\) is a rational
function of degree \(\leq 1\) in \(D^{-1}\) and \(C^{-1}\), (b) \(D\)
and \(C\) carry opposite signs, and (c) the traversal is Class A with
step \(\Delta\) derived from the Instantiation axis, without assuming
the value of \(\pi\) as input.}

Conjecture 8.4 is the key independence question. Establishing it would
show that the \(\pi/4\) grouping is structurally forced, not chosen post
hoc to match the Gregory--Leibniz series.

\begin{center}\rule{0.5\linewidth}{0.5pt}\end{center}

\subsection{9. Related work}\label{related-work}

\textbf{Series rearrangements.} The Leibniz series is conditionally
convergent, so rearrangements can alter its sum (Riemann rearrangement
theorem {[}Kno56{]}). The \(\pi/4\) branch uses \emph{grouped} (paired)
rearrangement, not a sign-pattern reordering, so its limit is standard.
Conjecture 8.4 would need to demonstrate that this specific pairing is
forced rather than selected.

\textbf{Continued fractions and metallic means.} The golden ratio is the
simplest nontrivial continued fraction, \(\varphi = [1;1,1,1,\ldots]\),
and the worst-approximable irrational {[}HW79{]}. The \(\varphi\) branch
recovers this via a Fibonacci-structured recurrence; the \(D\)--\(C\)
swap rule repackages the classical Fibonacci growth without competing
with continued-fraction theory {[}Kos01{]}.

\textbf{Newton--Babylonian method.} The \(\sqrt{2}\) branch is the
classical Heron iteration for \(x^2 = 2\) {[}Hea21, BF15{]}. The
tholonic framing adds a role interpretation (\(D\) as target bound,
\(C\) as averaging synthesizer) without altering the proof or the
convergence order.

\textbf{Classical series identities.} Standard analysis establishes
\(\sum_{k=0}^\infty 1/k! = e\) {[}Rud76{]} and
\(\sum_{k=1}^\infty (-1)^{k+1}/k = \ln 2\) {[}Kno56{]}. The ladder
presents these as Class B and C branches in a unified format; no claim
is made beyond this repackaging.

\textbf{Structural unification frameworks.} Several authors have
organized multiple classical constants under a single formalism; the
study of \emph{metallic means} extends \(\varphi\) to a parametric
family of continued fractions, a family the ladder realizes internally
(Proposition 6.6(ii)). We are unaware of a framework that identifies all
five of \(\pi/4\), \(\varphi\), \(e\), \(\sqrt{2}\), \(\ln 2\) as
branches of a common three-variable recurrence with consistent role
assignments. The closest conceptual territory is the general study of
fixed-point iterations and series convergence, but neither addresses the
triadic role separation or the seed-partition theorem.

\textbf{Leibniz and binary structure.} Leibniz, discoverer of the
\(\pi/4\) series, developed binary arithmetic simultaneously and drew
explicit connections to the I Ching {[}Lei03, Swe03{]}. The tholonic
framework formalizes a version of this intuition: binary state space
generates the minimum non-trivial simplex, the simplex induces the
triad, and the triad produces the five classical limits. The
interpretive layer is separable from the proofs.

\begin{center}\rule{0.5\linewidth}{0.5pt}\end{center}

\subsection{10. Discussion}\label{discussion}

\textbf{Organizational vs predictive claims.} The ladder framework as
proved is \emph{organizational}: a unified format and role vocabulary
from which five known limits descend. This is a genuine contribution
(unification with structural theorems attached) but weaker than a
\emph{predictive} claim of the form ``the triad axioms force these five
constants and no others at low seed complexity.'' The path from
organizational to predictive runs through the Conjectures of §8.

\textbf{Post hoc flexibility risk.} The standard referee objection to a
framework paper is: ``you chose seeds and rules after knowing the
targets.'' The structural theorems partially defuse this. Diagonal
invariance and swap symmetry are consequences of the Class A form
independent of any targeted limit. The seed partition (Proposition 6.1)
is a sharp combinatorial observation: a framework tuned post hoc would
not automatically produce a clean \(\{1,3,5\}\) vs \(\{0,1,2\}\)
partition. The perfect-square denominator identity (Proposition 6.4) is
an algebraic corollary of the step structure that a post hoc framework
would need to ``accidentally'' reproduce.

The remaining vulnerability is Conjecture 8.4. Until non-circularity is
established for the \(\pi/4\) grouping, one cannot rule out that it is
simply one among infinitely many valid Leibniz rearrangements that
happen to fit the Class A format.

\textbf{Convergence rates as a structural signal.} The two \(O(1/k)\)
branches are the Class A branch (\(\pi/4\)) and a Class C branch
(\(\ln 2\)); the other three reach machine precision within tens of
iterations (\(e\) super-exponentially, \(\sqrt{2}\) quadratically,
\(\varphi\) linearly at geometric ratio \(\varphi^{-2}\)). Note that the
per-step ordering places Class C \(\sqrt{2}\) (quadratic) above Class B
\(\varphi\) (linear), so the correlation between traversal class and
convergence order is imperfect; characterizing exactly when class
predicts rate is an open regularity whose resolution would add
quantitative predictive content to the classification of §4.1.

\textbf{The infinite landscape.} The five branches documented here are
recognizable loci in an infinite space of admissible seed--traversal
combinations. Proposition 6.6 exhibits two parametric families
explicitly: the square roots \(\sqrt{a}\) and the metallic means
\(\sigma_m\). Every distinct pair of seeds and traversal dynamics
defines a new path; every convergent path defines a new limit. The
scientific question is not whether other constants exist in this
landscape (they obviously do) but which of the infinite admissible paths
yield limits of independent mathematical significance, and whether that
set is finite and well-characterized. The Conjectures of §8 make this
question precise.

\begin{center}\rule{0.5\linewidth}{0.5pt}\end{center}

\subsection{11. Conclusion}\label{conclusion}

We have presented a three-variable recurrence family \((N, D, C)\), the
tholonic ladder, with five documented branches converging to \(\pi/4\),
\(\varphi\), \(e\), \(\sqrt{2}\), and \(\ln 2\). The limits are
classical; the contribution lies in the structural theorems: a seed
partition cleanly separating the \(\pi/4\) branch, diagonal invariance
and swap symmetry for the Class A template, a perfect-square denominator
identity linking \(\pi\) to the step structure, an offset-consistency
law for shifted initializations, and two parametric families (square
roots and metallic means) embedding the \(\sqrt{2}\) and \(\varphi\)
branches. Together these distinguish the five branches from an arbitrary
selection of known limits.

The central claim (\textbf{same recursive triadic geometry, different
traversal dynamics, different limits}) is supported by these proved
contrasts. The open work, stated as four numbered conjectures, concerns
the necessity of the seed assignments and traversal rules: whether they
are forced by a small set of axioms (predictive framework) or represent
a human-chosen sample from a larger admissible space (organizational
taxonomy). The structural evidence points toward the former; the proofs
are not yet in hand.

\begin{center}\rule{0.5\linewidth}{0.5pt}\end{center}

\subsection{Appendix A. Reference
implementation}\label{appendix-a.-reference-implementation}

The following Python code implements all five branches. No bare numeric
literals appear inside any recursion loop; only \(N\), \(D\), and \(C\)
(and the iteration index) appear in the update bodies.

\begin{Shaded}
\begin{Highlighting}[]
\CommentTok{\#!/usr/bin/env python3}

\KeywordTok{def}\NormalTok{ \_compute\_pi\_over\_4(max\_iter: }\BuiltInTok{int}\NormalTok{, istep: }\BuiltInTok{int} \OperatorTok{=} \DecValTok{2}\NormalTok{) }\OperatorTok{{-}\textgreater{}} \BuiltInTok{float}\NormalTok{:}
\NormalTok{    N }\OperatorTok{=} \FloatTok{1.0}
\NormalTok{    D }\OperatorTok{=} \FloatTok{3.0}
\NormalTok{    C }\OperatorTok{=} \FloatTok{5.0}
\NormalTok{    d\_step }\OperatorTok{=}\NormalTok{ istep }\OperatorTok{**} \DecValTok{2}      \CommentTok{\# = 4}
\NormalTok{    c\_step }\OperatorTok{=}\NormalTok{ istep }\OperatorTok{*} \DecValTok{2}       \CommentTok{\# = 4}
    \ControlFlowTok{for}\NormalTok{ \_ }\KeywordTok{in} \BuiltInTok{range}\NormalTok{(max\_iter):}
\NormalTok{        N }\OperatorTok{=}\NormalTok{ N }\OperatorTok{{-}}\NormalTok{ (}\FloatTok{1.0} \OperatorTok{/}\NormalTok{ D) }\OperatorTok{+}\NormalTok{ (}\FloatTok{1.0} \OperatorTok{/}\NormalTok{ C)}
\NormalTok{        D }\OperatorTok{+=}\NormalTok{ d\_step}
\NormalTok{        C }\OperatorTok{+=}\NormalTok{ c\_step}
    \ControlFlowTok{return}\NormalTok{ N}


\KeywordTok{def}\NormalTok{ \_compute\_phi(max\_iter: }\BuiltInTok{int}\NormalTok{) }\OperatorTok{{-}\textgreater{}} \BuiltInTok{float}\NormalTok{:}
\NormalTok{    N }\OperatorTok{=} \FloatTok{1.0}
\NormalTok{    D }\OperatorTok{=} \DecValTok{1}   \CommentTok{\# exact integer Fibonacci terms (floats overflow near k = 1476)}
\NormalTok{    C }\OperatorTok{=} \DecValTok{2}
\NormalTok{    base }\OperatorTok{=}\NormalTok{ N}
    \ControlFlowTok{for}\NormalTok{ \_ }\KeywordTok{in} \BuiltInTok{range}\NormalTok{(max\_iter):}
\NormalTok{        N }\OperatorTok{=}\NormalTok{ base }\OperatorTok{+}\NormalTok{ (D }\OperatorTok{/}\NormalTok{ C)}
\NormalTok{        D, C }\OperatorTok{=}\NormalTok{ C, (C }\OperatorTok{+}\NormalTok{ D)}
    \ControlFlowTok{return}\NormalTok{ N}


\KeywordTok{def}\NormalTok{ \_compute\_e(max\_iter: }\BuiltInTok{int}\NormalTok{) }\OperatorTok{{-}\textgreater{}} \BuiltInTok{float}\NormalTok{:}
\NormalTok{    N }\OperatorTok{=} \FloatTok{0.0}
\NormalTok{    D }\OperatorTok{=} \FloatTok{1.0}
\NormalTok{    C }\OperatorTok{=} \FloatTok{1.0}
    \ControlFlowTok{for}\NormalTok{ count }\KeywordTok{in} \BuiltInTok{range}\NormalTok{(max\_iter):}
\NormalTok{        N }\OperatorTok{=}\NormalTok{ N }\OperatorTok{+}\NormalTok{ (D }\OperatorTok{/}\NormalTok{ C)}
\NormalTok{        C }\OperatorTok{*=}\NormalTok{ count }\OperatorTok{+} \DecValTok{1} \ControlFlowTok{if}\NormalTok{ count }\OperatorTok{\textgreater{}} \DecValTok{0} \ControlFlowTok{else} \DecValTok{1}
    \ControlFlowTok{return}\NormalTok{ N}


\KeywordTok{def}\NormalTok{ \_compute\_sqrt2(max\_iter: }\BuiltInTok{int}\NormalTok{) }\OperatorTok{{-}\textgreater{}} \BuiltInTok{float}\NormalTok{:}
\NormalTok{    N }\OperatorTok{=} \FloatTok{1.0}
\NormalTok{    D }\OperatorTok{=} \FloatTok{2.0}
\NormalTok{    C }\OperatorTok{=} \FloatTok{2.0}
    \ControlFlowTok{for}\NormalTok{ \_ }\KeywordTok{in} \BuiltInTok{range}\NormalTok{(max\_iter):}
\NormalTok{        N }\OperatorTok{=}\NormalTok{ (N }\OperatorTok{+}\NormalTok{ (D }\OperatorTok{/}\NormalTok{ N)) }\OperatorTok{/}\NormalTok{ C}
    \ControlFlowTok{return}\NormalTok{ N}


\KeywordTok{def}\NormalTok{ \_compute\_ln2(max\_iter: }\BuiltInTok{int}\NormalTok{) }\OperatorTok{{-}\textgreater{}} \BuiltInTok{float}\NormalTok{:}
\NormalTok{    N }\OperatorTok{=} \FloatTok{0.0}
\NormalTok{    D }\OperatorTok{=} \FloatTok{1.0}
\NormalTok{    C }\OperatorTok{=} \FloatTok{1.0}
    \ControlFlowTok{for}\NormalTok{ count }\KeywordTok{in} \BuiltInTok{range}\NormalTok{(max\_iter):}
\NormalTok{        term }\OperatorTok{=}\NormalTok{ D }\OperatorTok{/}\NormalTok{ (count }\OperatorTok{+}\NormalTok{ C)}
\NormalTok{        N }\OperatorTok{=}\NormalTok{ N }\OperatorTok{+}\NormalTok{ term }\ControlFlowTok{if}\NormalTok{ count }\OperatorTok{\%} \DecValTok{2} \OperatorTok{==} \DecValTok{0} \ControlFlowTok{else}\NormalTok{ N }\OperatorTok{{-}}\NormalTok{ term}
    \ControlFlowTok{return}\NormalTok{ N}


\KeywordTok{def}\NormalTok{ compute\_tholonic\_constant(constant\_type}\OperatorTok{=}\StringTok{"pi/4"}\NormalTok{, max\_iter}\OperatorTok{=}\DecValTok{100000}\NormalTok{):}
\NormalTok{    calculators }\OperatorTok{=}\NormalTok{ \{}
        \StringTok{"pi/4"}\NormalTok{:  \_compute\_pi\_over\_4,}
        \StringTok{"phi"}\NormalTok{:   \_compute\_phi,}
        \StringTok{"e"}\NormalTok{:     \_compute\_e,}
        \StringTok{"sqrt2"}\NormalTok{: \_compute\_sqrt2,}
        \StringTok{"ln2"}\NormalTok{:   \_compute\_ln2,}
\NormalTok{    \}}
\NormalTok{    calc }\OperatorTok{=}\NormalTok{ calculators.get(constant\_type)}
    \ControlFlowTok{if}\NormalTok{ calc }\KeywordTok{is} \VariableTok{None}\NormalTok{:}
        \ControlFlowTok{raise} \PreprocessorTok{ValueError}\NormalTok{(}\SpecialStringTok{f"Unknown constant\_type: }\SpecialCharTok{\{}\NormalTok{constant\_type}\SpecialCharTok{\}}\SpecialStringTok{"}\NormalTok{)}
    \ControlFlowTok{return}\NormalTok{ calc(max\_iter)}


\ControlFlowTok{if} \VariableTok{\_\_name\_\_} \OperatorTok{==} \StringTok{"\_\_main\_\_"}\NormalTok{:}
    \ControlFlowTok{for}\NormalTok{ ct }\KeywordTok{in}\NormalTok{ [}\StringTok{"pi/4"}\NormalTok{, }\StringTok{"phi"}\NormalTok{, }\StringTok{"e"}\NormalTok{, }\StringTok{"sqrt2"}\NormalTok{, }\StringTok{"ln2"}\NormalTok{]:}
        \BuiltInTok{print}\NormalTok{(}\SpecialStringTok{f"}\SpecialCharTok{\{}\NormalTok{ct}\SpecialCharTok{:6s\}}\SpecialStringTok{ = }\SpecialCharTok{\{}\NormalTok{compute\_tholonic\_constant(ct)}\SpecialCharTok{\}}\SpecialStringTok{"}\NormalTok{)}
\end{Highlighting}
\end{Shaded}

\textbf{Output at \(10^5\) iterations:}

\begin{verbatim}
pi/4   = 0.7853994133911905
phi    = 1.618033988749895
e      = 2.7182818284590455
sqrt2  = 1.414213562373095
ln2    = 0.6931421805849816
\end{verbatim}

\begin{center}\rule{0.5\linewidth}{0.5pt}\end{center}

\subsection{Appendix B. Summary table with convergence
rates}\label{appendix-b.-summary-table-with-convergence-rates}

\begin{longtable}[]{@{}
  >{\raggedright\arraybackslash}p{(\linewidth - 10\tabcolsep) * \real{0.1159}}
  >{\raggedright\arraybackslash}p{(\linewidth - 10\tabcolsep) * \real{0.2464}}
  >{\raggedright\arraybackslash}p{(\linewidth - 10\tabcolsep) * \real{0.1449}}
  >{\raggedright\arraybackslash}p{(\linewidth - 10\tabcolsep) * \real{0.1014}}
  >{\raggedright\arraybackslash}p{(\linewidth - 10\tabcolsep) * \real{0.3043}}
  >{\raggedright\arraybackslash}p{(\linewidth - 10\tabcolsep) * \real{0.0870}}@{}}
\toprule\noalign{}
\begin{minipage}[b]{\linewidth}\raggedright
Target
\end{minipage} & \begin{minipage}[b]{\linewidth}\raggedright
\((N_0,D_0,C_0)\)
\end{minipage} & \begin{minipage}[b]{\linewidth}\raggedright
Seed set
\end{minipage} & \begin{minipage}[b]{\linewidth}\raggedright
Class
\end{minipage} & \begin{minipage}[b]{\linewidth}\raggedright
Classical mechanism
\end{minipage} & \begin{minipage}[b]{\linewidth}\raggedright
Rate
\end{minipage} \\
\midrule\noalign{}
\endhead
\bottomrule\noalign{}
\endlastfoot
\(\pi/4\) & \((1,3,5)\) & \(\{1,3,5\}\) & A & Leibniz paired grouping &
\(O(1/K)\) \\
\(\varphi\) & \((1,1,2)\) & \(\{1,2\}\) & B & Fibonacci ratio
convergence & linear, rate \(\varphi^{-2}\) \\
\(e\) & \((0,1,1)\) & \(\{0,1\}\) & B & \(\sum_{k=0}^\infty 1/k!\) &
\(O(1/K!)\) \\
\(\sqrt{2}\) & \((1,2,2)\) & \(\{1,2\}\) & C & Newton--Babylonian &
quadratic \\
\(\ln 2\) & \((0,1,1)\) & \(\{0,1\}\) & C & alternating harmonic &
\(O(1/K)\) \\
\end{longtable}

Offset behavior under non-zero \(N_0\) initializations: see Proposition
6.5.

\begin{center}\rule{0.5\linewidth}{0.5pt}\end{center}

\subsection{References}\label{references}

{[}Lei82{]} G. W. Leibniz, ``De vera proportione circuli ad quadratum
circumscriptum in numeris rationalibus,'' \emph{Acta Eruditorum}, 1682.

{[}Lei03{]} G. W. Leibniz, ``Explication de l'Arithmétique Binaire,''
\emph{Mémoires de l'Académie Royale des Sciences}, 1703.

{[}Roy11{]} R. Roy, \emph{Sources in the Development of Mathematics:
Infinite Series and Products from the Fifteenth to the Twenty-first
Century}, Cambridge University Press, 2011.

{[}Kno56{]} K. Knopp, \emph{Infinite Sequences and Series}, Dover, 1956.

{[}HW79{]} G. H. Hardy and E. M. Wright, \emph{An Introduction to the
Theory of Numbers}, 5th ed., Oxford University Press, 1979.

{[}Kos01{]} T. Koshy, \emph{Fibonacci and Lucas Numbers with
Applications}, Wiley, 2001.

{[}Hea21{]} T. L. Heath, \emph{A History of Greek Mathematics}, Vol. 1,
Oxford University Press, 1921.

{[}BF15{]} R. L. Burden and J. D. Faires, \emph{Numerical Analysis},
10th ed., Cengage Learning, 2015.

{[}Rud76{]} W. Rudin, \emph{Principles of Mathematical Analysis}, 3rd
ed., McGraw-Hill, 1976.

{[}Mil24{]} J. W. Milton, \emph{Tholonia: The Existential Mechanics of
Awareness}, self-published, 2024. Available at
https://tholonia.github.io/the-book/.

{[}Swe03{]} F. J. Swetz, ``Leibniz, the Yijing, and the Religious
Conversion of the Chinese,'' \emph{Mathematics Magazine} 76(4), 2003,
pp.~276--291.

\end{document}
